\section{Muon Fake Factor Measurement}%
\label{sec:bkg-fakes-muon}

Fake muons arise from in-flight decays of mesons inside jets with
the semi-leptonic decay of B meson as one of the largest contributions.
Combinations of prompt and fake muons can contribute to the measurement region
if the fake muons satisfy tight isolation and impact parameter requirements.
The muon fake contribution to the measurement region is estimated as described
in \cref{sec:bkg-fakefactor}, which involves a data driven estimate
of the fake factor.


\subsection{Muon Fake Factor Measurement}

Fake-factors are measured requiring a jet to be associated to a muon, thereby
targeting a fake factor measurement in di-jet events. Basing on the selection defined
for electrons in \cref{tab:FakeRegionDefinitions}, events are selected
by requiring the highest \pt light jet in the event to have an azimuthal angular distance
greater than \( \Delta \phi > 2.7 \) to the muon. This topology is qualitatively illustrated
in \cref{fig:bkg-dijet}.

\begin{figure}[htbp]
  \begin{center}
    \includegraphics[width=0.6\textwidth]{fakes/dijet}
    \caption{Qualitative sketch of a di-jet event selected for the measurement
      of the muon fake-factors.}%
    \label{fig:bkg-dijet}
  \end{center}
\end{figure}

The selection of di-jet events allows one to achieve high statistics for
measuring fake-factors, while also allowing to address systematic uncertainties
using the properties of the recoiling tag jet.

The muon and the leading=\pt jet in the event have to be back-to-back in the transverse
plane, i.e. \( \Delta \phi(\mu, \mathrm{jet}) > 2.7 \) and have \pt threshold
of \( \pt(\mu) > \qty{30}{\GeV} \) and \( \pt(\mathrm{jet}) > \qty{35}{\GeV} \).
The selection criteria of the muon fake measurement region are summarised in \cref{tab:FakeRegionDefinitions}.

The distributions of muon \pt and \eta after all selection requirements
are applied, are reported in \cref{fig:Fakes_Muons-distributions}, where
tight and loose muons are tight and loose, respectively, as defined in
\cref{tab:definitions-muons}. The \pt dependence on fake factors is shown
in \cref{fig:Fakes_Muons}.

\begin{figure}[htbp]
  \centering
  \includegraphics[width=0.4\textwidth]{fakes/fake_factors_m_nominal}
  \caption{Muon fake factors measured as a function of \pt.}%
  \label{fig:Fakes_Muons}
\end{figure}

\begin{figure}[htbp]
  \centering
  \subfloat[]{
    \includegraphics[page=3, width=0.4\textwidth]{fakes/Fakes_Muons_m_nominal}%
    \label{fig:Fakes_Muons-pt-loose}
  }
  \subfloat[]{
    \includegraphics[page=1, width=0.4\textwidth]{fakes/Fakes_Muons_m_nominal}%
    \label{fig:Fakes_Muons-pt-tight}
  }
  \\
  \subfloat[]{
    \includegraphics[page=6, width=0.4\textwidth]{fakes/Fakes_Muons_m_nominal}%
    \label{fig:Fakes_Muons-eta-loose}
  }
  \subfloat[]{
    \includegraphics[page=5, width=0.4\textwidth]{fakes/Fakes_Muons_m_nominal}%
    \label{fig:Fakes_Muons-eta-tight}
  }

  \caption{Fakes-enriched regions in the nominal selection.
    \protect\subref{fig:Fakes_Muons-pt-loose} \pt distribution of loose muons in the fakes-enriched region,
    \protect\subref{fig:Fakes_Muons-pt-tight} \pt distribution of tight muons in the fakes-enriched region,
    \protect\subref{fig:Fakes_Muons-eta-loose} \eta distribution of loose muons in the fakes-enriched region and
    \protect\subref{fig:Fakes_Muons-eta-tight} \eta distribution of tight muons in the fakes-enriched region.
    All the distributions show data events and the prompt MC component subtracted from data, to ensure a fake dominated region.
  }%
  \label{fig:Fakes_Muons-distributions}
\end{figure}


\subsection{Systematic Uncertainties on Muon Fake Factors}%
\label{sec:bkg-fakes-muon-sys}

Similar to electrons all Monte Carlo samples were varied by \qty{10}{\%}
to account for cross-section and luminosity uncertainties as well as the
modelling of the prompt MC used in the subtraction procedure.

Additionally the effect of the \( \met + \mtW \) cut is tested by varying it
to \qty{50}{\GeV} and \qty{70}{\GeV}. The effect of this variation on the fake
factors is not negligible.

The effect on the number of \bjets was tested by varying \( N_{\bjets}\) to
include at least 2 \bjets. This effect is not negligible and can be quite large
in some bins in the electron \pt distribution. 

Additional variations were tested, including the variation of the transverse
momentum of the leading jet, with thresholds of \( \pt(j1) > \qty{40}{\GeV} \)
and \qty{30}{\GeV}, \( \Delta \phi(\mu, \mathrm{jet})\) > 2.6 and 2.8,
and the effect of omitting the requirement on light jets.

The effect of each systematic alteration on the nominal measurement can be
seen on \crefrange{fig:Fakes_Muons-musys1}{fig:Fakes_Muons-musys3}.
The total uncertainty is estimated by comparing the statistical uncertainty on
the nominal measurement with the maximum deviation between the nominal
fake-factor and each systematically altered measurement. The largest signed
deviation is taken to be the total uncertainty for lower bounds and upper bounds
in turn.

\begin{figure}[htbp]
  \centering
  \subfloat[]{
    \includegraphics[page=1, width=0.4\textwidth]{fakes/fake_factors_m_final}%
    \label{fig:Fakes_Muons-musys1}
  }
  \subfloat[]{
    \includegraphics[page=2, width=0.4\textwidth]{fakes/fake_factors_m_final}%
    \label{fig:Fakes_Muons-musys2}
  }
  \\
  \subfloat[]{
    \includegraphics[page=3, width=0.4\textwidth]{fakes/fake_factors_m_final}%
    \label{fig:Fakes_Muons-musys3}
  }

  \caption{The measured fake factor for each systematics variation.
    \protect\subref{fig:Fakes_Muons-musys1} fake factor variations for the first \eta bin,
    \protect\subref{fig:Fakes_Muons-musys2} fake factor variations for the second \eta bin and
    \protect\subref{fig:Fakes_Muons-musys3} fake factor variations for the third \eta bin.
  }% 
  \label{fig:Fakes_Muons-musys}
\end{figure}


\subsection{Validation of Fake-factors for Muons}%
\label{subsec:bkg-fakes-muon-validation}

The measured fake factors were validated by performing the direct closure test
in the fakes-enriched region. The same event selection was used as for the
estimation, but no special binning has been applied. Fakes have been estimated
from data using the measured fake factors. The reconstructed object selection
was the same as the nominal event selection (\cref{ch:definitions}).

Prompt Monte Carlo processes that were used are \Wjets, \Zjets, \ttbar and
single top, diboson, and rare top decays. The remaining events
are described well by the fake factor method as seen on
\cref{fig:Fakes_ClosureMuons}.
The only considered uncertainty in this region is the systematic uncertainty
of the fake factor.

\begin{figure}[htbp]
  \centering
  \subfloat[]{
    \includegraphics[page=3, width=0.4\textwidth]{fakes/Fakes_ClosureMuons_m_validation_sys}%
    \label{fig:Fakes_ClosureMuons-pt}
  }
  \subfloat[]{
    \includegraphics[page=7, width=0.4\textwidth]{fakes/Fakes_ClosureMuons_m_validation_sys}%
    \label{fig:Fakes_ClosureMuons-eta}
  }
  \\
  \subfloat[]{
    \includegraphics[page=9, width=0.4\textwidth]{fakes/Fakes_ClosureMuons_m_validation_sys}%
    \label{fig:Fakes_ClosureMuons-met}
  }
  \subfloat[]{
    \includegraphics[page=15, width=0.4\textwidth]{fakes/Fakes_ClosureMuons_m_validation_sys}%
    \label{fig:Fakes_ClosureMuons-mt}
  }

  \caption{The muon fake-factors closure.
    \protect\subref{fig:Fakes_ClosureMuons-pt} muon \pt distribution and
    \protect\subref{fig:Fakes_ClosureMuons-eta} muon \eta distribution,
    \protect\subref{fig:Fakes_ClosureMuons-met} \met distribution, and
    \protect\subref{fig:Fakes_ClosureMuons-mt} \mtW distribution.
    The dashed band shows the systematic uncertainty associated to the fake factor measurement.
  }%
  \label{fig:Fakes_ClosureMuons}
\end{figure}
